\section{Limitations} 
\label{sec: Limitations}

\subsection[Absolute CDA Interpretation]{Absolute $C_DA$ Interpretation}
\label{sec: Absolute CDA Interpretation}

At Jeddah the composite $2\hat\alpha/\rho = 1.485$~m$^2$ implies $C_DA = 1.448$~m$^2$, which exceeds the pre-2022 benchmark of $0.8$--$1.0$~m$^2$ sometimes cited in the literature but sits squarely within the range expected for the 2022+ ground-effect generation. Three lines of evidence show this is a genuine measurement rather than a systematic inflation. (One previously suspected inflation source---a driver-mass double-count---has been removed: the FIA $798$~kg minimum already includes the driver, and the corrected mass model is used throughout; see \S\ref{sec: Coast Down ODE Model}.)

\textbf{Gear-stratified falsification.} If engine braking were the primary cause, gear-8-only segments---where aerodynamic drag at high speed is $15$--$20\times$ larger than engine braking---should yield substantially lower $\alpha$. At Jeddah (13 segments; 9 in gear~8), the gear-8-only fit gives $\hat{\alpha} = 0.867$, within $1.1\%$ of the all-segment value of $0.876$. The residual ($-1.0\%$) is not statistically distinguishable from zero within the bootstrap 90\% CI ($\pm 3.1\%$), confirming that engine braking is not a material contributor to the composite.

\textbf{2022+ generation benchmark.} The ground-effect cars introduced in 2022 carry significantly more base drag than pre-2022 cars due to larger bodywork and sealed wheel-fairing geometry. Published open-source lap simulation estimates (see, e.g., vehicle dynamics community benchmarks) place $C_DA$ for 2022+ cars at Jeddah in the range $1.3$--$1.6$~m$^2$; this range is indicative rather than authoritative, as no primary source has been independently verified. The measured $1.45$~m$^2$ sits in the lower-middle of this band, consistent with Jeddah's medium-low downforce package. At the high end, Suzuka at $1.69$~m$^2$ is consistent with known high-downforce configurations for the 2022+ cars.

\textbf{Parameter stability and $\beta$ sensitivity.} The $v_0$-fitting result (\S\ref{sec: Pooled Fitting}), the near-invariance of $\hat{\alpha}$ to the $\beta$ prior (\S\ref{sec: Coast Down ODE Model}), and the gear stratification together confirm that $\hat{\alpha}$ is robust: it is stable to initial-condition assumptions, to the rolling-resistance prior, and to gear selection, all of which would be expected to shift the estimate if a missing force term were inflating it. To quantify the residual sensitivity, the Jeddah coast-down segments were re-fitted with $\beta$ swept across [80, 100, 120, 140, 160]~N (see \texttt{scripts/beta\_sensitivity.py}). The resulting $C_DA$ varied from $+1.2\%$ (at $\beta = 80$~N) to $-1.5\%$ (at $\beta = 160$~N) relative to the $\beta = 120$~N baseline, an absolute range of $1.427$--$1.467$~m$^2$. This confirms that a $\pm 33\%$ variation in $\beta$ shifts absolute $C_DA$ by less than $\pm 1.5\%$, which is small relative to the dominant $\mu$-systematic uncertainty in the polar (see \S\ref{sec: GPS Noise}).

Remaining $\text{DW} < 2$ after $v_0$ fitting (DW~$\approx 0.55$; improved from $0.35$ before $v_0$ fitting, see \S\ref{sec: Pooled Fitting}) reflects aerodynamic pitch and ride-height dynamics during coast-down---$C_DA$ varies slightly as the car pitches forward under deceleration---which is not captured by the constant-$C_DA$ model. This is a structural limitation of the coast-down method, not correctable without ride-height telemetry.

A practical consequence of DW $\approx 0.55$ is that the covariance matrix produced by \texttt{scipy.optimize.least\_squares} is anti-conservative: positive residual autocorrelation inflates the apparent goodness-of-fit and deflates the reported parameter standard errors. The quoted $\pm 0.004$~(N\,s$^2$/m$^2$) on $\alpha$ (Jeddah) and the resulting $\sim\!\pm 0.007$~m$^2$ on $C_DA$ should therefore be treated as lower bounds on the true fitting uncertainty. The circuit-level bootstrap CIs and the $13\times$ conservative ratio are first quoted in \S\ref{sec: Jeddah Validation}; the full justification follows. A segment-resampling bootstrap ($n=1000$ replicates) implemented in \texttt{src/uncertainty.py} provides a more honest characterisation: it resamples the available segments with replacement, re-runs the pooled fit on each replicate, and returns a percentile CI that captures segment-to-segment variation without relying on the anti-conservative covariance matrix. 

The bootstrap 90\% CI on $\hat\alpha$ at Jeddah spans $\pm 0.036$~N\,s$^2$/m$^2$ (see \texttt{scripts/recompute\_report\_numbers.py}), and should be used in preference to the Jacobian estimate for any downstream decision that depends on the absolute magnitude of $\hat\alpha$. One residual limitation of segment resampling is that segments within the same FP session share common systematic effects (track state, ambient temperature, session-level fuel load), so they are not fully exchangeable. If the dominant uncertainty is session-level rather than segment-level, the segment bootstrap understates the between-session variance.

To directly test this, a session-stratified bootstrap ($n=1000$ replicates; \texttt{bootstrap\_alpha\_ci\_session\_stratified} in \texttt{src/uncertainty.py}) was run at Jeddah. This resamples whole FP sessions with replacement rather than individual segments, producing a between-session CI. At Jeddah, 13 segments pass quality filters across two contributing sessions (FP1 yielded no passing segments; FP2: 3 segments, FP3: 10 segments). The session-stratified 90\% CI is $\pm 0.014$~N\,s$^2$/m$^2$---narrower than the segment-level CI of $\pm 0.036$, a ratio of $0.39\times$. This indicates that between-session $\hat\alpha$ variance at Jeddah is small relative to segment-to-segment noise: FP2 and FP3 give consistent $\hat\alpha$ values, suggesting track evolution across sessions was not a dominant source of variation at this venue. The caveat is that with only two effective sessions, the session-level bootstrap can produce only three distinct resample outcomes (FP2-only, FP3-only, or mixed), so the CI is coarsely discretised; the $0.39\times$ ratio should be read as indicative rather than precise. In short, the segment bootstrap CI of $\pm 0.036$ is a conservative bound at Jeddah; session-level systematics appear smaller than within-session segment noise at this circuit. 

This ratio is lower than the $\sim\!13\times$ observed at Monza in earlier analysis (re-evaluated using the corrected Jacobian covariance formula; see \S\ref{sec: Pooled Fitting}), reflecting less residual autocorrelation in the Jeddah segments (consistent with a lower-downforce, shorter-segment regime). 

Monza was excluded from the polar on data-quality grounds (only 3 valid qualifying push laps), but it represents a lower-downforce, longer-segment regime with worse autocorrelation than any polar circuit; using it as the conservative upper bound is therefore not circular, and errs on the side of wider $C_DA$ uncertainty. \begin{table}[h]
\centering
\caption{Bootstrap-to-Jacobian ratio used to scale CDA uncertainty in the polar Monte Carlo. The conservative $13\times$ from Monza is applied at all circuits.}
\label{tab: bootstrap scaling}
\begin{tabularx}{\linewidth}{l c c l}
\toprule
Circuit & {Jacobian $\sigma_\alpha$ (N\,s$^2$/m$^2$)} & {Bootstrap 90\% CI half-width (N\,s$^2$/m$^2$)} & Ratio \\
\midrule
Jeddah   & 0.004 & 0.027 & $6.8\times$ \\
Monza$^\dagger$ & 0.008 & $\sim\!0.104$ & $13\times$ (conservative upper bound) \\
\bottomrule
\end{tabularx}
\footnotesize{$^\dagger$Monza was excluded from the polar on data-quality grounds (only 3 valid qualifying push laps) but provides the conservative upper bound for the scaling.}
\end{table}

A conservative $13\times$ scaling (Table~\ref{tab: bootstrap scaling}) is applied in the polar Monte Carlo: the Jeddah-calibrated $6.8\times$ is the lower bound and the Monza $13\times$ is the upper bound; applying the upper bound errs toward wider rather than narrower confidence intervals, avoiding under-statement of $C_DA$ uncertainty at circuits where the ratio was not independently calibrated. The cross-circuit \emph{ratio} and $C_{rr}$ are less affected, as systematic autocorrelation biases cancel when comparing fits run under identical methodology.

The per-circuit $C_DA$ uncertainty (Eq.~\ref{eq: cda_uncertainty}, \S\ref{sec: Deriving CDA}) receives contributions from both $\hat\alpha$ and $C_LA$. At Jeddah ($\sigma_{\hat\alpha} = 0.004$~N\,s$^2$/m$^2$, $\rho = 1.180$~kg/m$^3$, $\sigma_{C_LA} = 0.42$~m$^2$):
\begin{equation}
    \sigma_{C_DA} \approx \sqrt{(0.007)^2 + (0.006)^2} \approx 0.009\;\text{m}^2
\end{equation}
with the $C_{rr}\cdot\sigma_{C_LA}$ term ($\approx 0.006$~m$^2$) contributing comparably to the $\alpha$-fitting term ($\approx 0.007$~m$^2$). Because the Jacobian $\sigma_{\hat\alpha}$ is anti-conservative (DW~$\approx 0.55$; see above), $\pm 0.009$~m$^2$ is a lower bound on the true $C_DA$ uncertainty, and the $\pm$ values in Table~\ref{tab: five circuit results} should be read accordingly.

The cross-circuit composite ratios (e.g.\ Suzuka/Jeddah $= 1.18\times$) are unaffected by these considerations---the residual autocorrelation is common-mode and cancels between circuits run under identical methodology---and relative values remain the method's most reliable output.

\textbf{Free-practice fuel load.} The coast-down segments use FP1, FP2, and FP3 data. The per-lap mass model (\S\ref{sec: Coast Down ODE Model}) resets the initial fuel load to 95~kg at the start of each session independently. In practice, teams vary fuel loads substantially across free-practice sessions (short-run low-fuel laps, long race-simulation stints on high fuel), so the actual fuel load at the start of any given session may differ from 95~kg. A 20~kg error in the assumed initial load represents a $\sim\!2.3\%$ error in mass at mid-stint (at $m \approx 855$~kg), which propagates linearly to an equal fractional error in $\hat\alpha$. This effect adds unquantified scatter across FP sessions that is absorbed by the bootstrap CI but not separately attributed.

\textbf{Cross-session mass inconsistency.} The coast-down ODE (\S\ref{sec: Coast Down ODE Model}) uses a per-lap FP mass model (ranging from $\sim\!798$~kg at end-of-stint to $\sim\!893$~kg at full fuel load), while the $C_LA$ estimate (\S\ref{sec: ClLA from GPS Lateral Acceleration}) uses qualifying mass ($803$~kg). Both $\hat{\alpha}$ and $C_LA$ scale approximately linearly with their respective assumed masses, so when combined in Eq.~(\ref{eq: cda_identity}), a difference between the mass assumptions introduces a systematic offset in $C_DA$. At Jeddah, the mean FP fitting mass of $\sim\!855$~kg versus qualifying mass of $803$~kg is a $\sim\!6.5\%$ differential; since $C_{rr}\cdot C_LA \approx 0.037$~m$^2$ is only $\sim\!2.5\%$ of $C_DA$, the dominant effect is an $\sim\!6\%$ inflation of the $2\hat{\alpha}/\rho$ term relative to what a single consistent mass would give. This systematic can be bounded explicitly: $2\hat{\alpha}/\rho \approx 1.485$~m$^2$ at Jeddah, so a 6.5\% inflation from the FP-vs-qualifying mass mismatch corresponds to an absolute offset of $\sim\!0.097$~m$^2$ --- comparable to the $0.11$~m$^2$ spread between Yas Marina and Jeddah in Table~\ref{tab: five circuit results}. The effect is absorbed by the bootstrap CI but is not separately attributed and constitutes an unquantified floor on the absolute $C_DA$ values. Cross-circuit ratios are less affected, as the same FP mass model is applied at every circuit.

\textbf{Track gradient.} The coast-down ODE omits the gravitational resistance term $mg\sin\alpha$ present in full longitudinal vehicle models. At circuits with significant elevation change---notably Spa-Francorchamps and Suzuka---this is a potential source of segment-level bias: a coast-down occurring on a downhill gradient would show lower-than-expected deceleration, causing the fitter to underestimate $\hat\alpha$, while an uphill gradient would inflate it. In practice, the segment filter (entry speed $>180$~km/h, throttle $<5\%$, brake $= 0$, duration $>0.5$~s) selects events that predominantly occur on long high-speed straights, which are the flattest sections of any circuit; the pronounced elevation changes at Spa (Eau Rouge/Raidillon) and Suzuka (Degner) occur mid-corner at speeds excluded by the filter. Adding a gradient correction would require reliable per-sample elevation data: FastF1 exposes GPS altitude at 10~Hz, but double-differencing altitude at this rate introduces noise comparable to the gradient signal itself, making a correction unreliable without a surveyed elevation map. The gradient effect is therefore not modelled; it contributes unquantified scatter to circuits with non-flat straights but is expected to be small relative to the bootstrap CI given the segment selection criteria.

\textbf{Drivetrain rotational inertia.} The coast-down ODE models inertia as pure translational mass $m$. A more complete longitudinal model adds the effective translational mass of rotating drivetrain components:
\begin{equation}
    m_\text{eff} = m + \frac{n^2 I_e + 4\,I_w}{r^2}
\end{equation}
where $n$ is the overall gear ratio (engine speed / wheel speed), $I_e$ is engine rotational inertia, $I_w$ is per-wheel inertia (including brake disc), and $r$ is tyre rolling radius. For an F1 car in gear~8, a rough estimate using $I_w \approx 1$~kg\,m$^2$ per corner and $r \approx 0.33$~m gives $4I_w/r^2 \approx 37$~kg; the engine contribution is smaller but non-negligible. The total rotational inertia correction is on the order of $3$--$5\%$ of vehicle mass ($\sim\!25$--$45$~kg on an $855$~kg car), propagating linearly to an equal fractional inflation of $\hat\alpha$ and $C_DA$. However, the correction is strictly common-mode: the same engine and drivetrain are used at every circuit, and most segments are in gear~8 (\S\ref{sec: Jeddah Validation}). A common-mode inflation cancels in the cross-circuit ratios and does not affect the polar slope. Proprietary F1 drivetrain inertia values are not publicly available, so the correction cannot be applied quantitatively; it contributes an unquantified systematic offset to absolute $C_DA$ of similar order to the cross-session mass inconsistency discussed above.

\textbf{Speed-dependent rolling resistance.} The model uses a speed-independent rolling resistance $\beta = C_{rr}\cdot mg$, but at high speed ($>200$~km/h) F1 tyres exhibit a velocity-dependent viscous loss component that adds a retarding force proportional to $v$ or $v^2$. This component is indistinguishable from aerodynamic drag during coast-down and is therefore absorbed into $\hat\alpha$, causing a slight inflation of the composite $2\hat\alpha/\rho$ and hence the absolute $C_DA$ floor. Because the effect is common-mode across all circuits (the same tyre construction and operating speed range is used everywhere), it cancels in the cross-circuit ratios and does not affect the polar slope; its contribution to the absolute $C_DA$ offset is unquantified.

\subsection{DRS Inseparable From Slipstream and Race-Lap Confounds}
\label{sec: DRS Inseparable From Slipstream}

The gap-stratification analysis confirms that reducing slipstream reduces the apparent delta but does not eliminate it. Beyond slipstream, three additional first-order confounds affect the race-lap DRS energy balance and are not correctable from available data:

\textbf{Fuel load.} Race laps span the full stint, during which mass decreases by approximately $1.8$~kg/lap. A $50$~kg fuel burn over a 28-lap stint reduces mass by $\sim 6\%$, shifting trap speed by $\sim 0.3$~km/h --- small relative to the $11.9$~km/h observed delta, but not negligible relative to the $90\%$ CI width. The per-lap mass model (\S\ref{sec: Coast Down ODE Model}) is applied in the energy balance, which accounts for most of this effect.

\textbf{Tyre state.} Rolling resistance varies with tyre temperature, compound wear, and graining. DRS-open laps are concentrated in the early-to-mid stint (when drivers attack for position), while DRS-closed laps span the full distribution. A systematic offset in tyre $C_{rr}$ between these groups would bias $\Delta C_DA$.

\textbf{ERS deployment strategy.} MGU-K deployment on the main straight varies lap-by-lap with battery state-of-charge and team strategy. DRS-open laps may receive different MGU-K boosts than DRS-closed laps, directly confounding the trap speed comparison. This is the hardest confound to bound without internal ERS telemetry.

Together, these confounds mean that the $90\%$ CI on $\Delta C_DA$ ($-0.26$ to $+0.05$~m$^2$) is an optimistic lower bound on the true uncertainty. A clean DRS measurement from race data is not achievable with public telemetry; qualifying or free-practice track-mapping data (where DRS is active, ERS deployment is consistent lap-to-lap, and following cars are rare) would be required.

\subsection[CLA Depends on Assumed mu]{$C_LA$ Depends on Assumed $\mu$}
\label{sec: CLA Depends on Assumed mu}

The GPS lat-g method requires an assumed tyre friction coefficient. The sensitivity (roughly $C_LA \propto 1/\mu$) means a 10\% error in $\mu$ produces a 10\% error in $C_LA$. Published values for Pirelli compounds at operating temperature range from approximately $1.5$ to $2.0$. For the Jeddah qualifying dataset ($C_LA = 2.66$~m$^2$ at $\mu = 1.8$), this places $C_LA$ in the range $2.4$--$3.2$~m$^2$; at Suzuka ($4.33$~m$^2$) the band is $3.9$--$5.2$~m$^2$. Two distinct effects must be kept separate. \emph{Uniform $\mu$ error} (the same offset at every circuit): because the same $\mu = 1.8$ is assumed everywhere, this scales all $C_LA$ values by the same factor and shifts the polar intercept without changing the slope. \emph{Circuit-to-circuit $\mu$ variation} (different compounds, track temperatures, tyre ages): this moves each circuit's $C_LA$ by a different amount, rotating the polar and directly biasing the slope. The $\mu$ uncertainty therefore contributes to \emph{absolute} $C_LA$ values via the first effect, and to the \emph{polar slope} via the second. This circuit-to-circuit $\mu$ variation is unquantified and constitutes the dominant residual systematic in the polar slope estimate.

\subsection{GPS Noise and Precision of the $C_LA$ Estimate}
\label{sec: GPS Noise}

Computing lateral acceleration by double-differencing 10~Hz GPS position data introduces substantial noise. A 5-sample rolling mean (0.5~s window) is insufficient to suppress the resulting scatter, as seen in the large per-lap error bars in Figure~3.2 (individual $\pm 1\sigma$ intervals spanning $1$--$2$~m$^2$). The estimate relies on the median over $N$ laps to converge; its precision is characterised by the standard error of the median:
\begin{equation}
    \text{SE}_\text{median} \approx 1.253 \cdot \frac{\sigma_\text{lap}}{\sqrt{N}}
\end{equation}
Here $\sigma_\text{lap}$ is the lap-to-lap scatter of the per-lap median estimates. Across the five polar circuits, the statistical SE of the median ranges from $0.13$~m$^2$ (Jeddah and Spa, where qualifying push laps produce consistent estimates) to $0.19$~m$^2$ (Silverstone). In all cases the dominant remaining uncertainty is the assumed $\mu$ (\S\ref{sec: CLA Depends on Assumed mu}), which contributes $0.40$--$0.65$~m$^2$ per circuit and is the binding constraint on the polar uncertainty.

\subsection{Driver Independence}

All five circuits in this analysis use driver~14 (Alonso, AMR24), eliminating the cross-driver confound present in earlier versions that compared driver~14 at Monza with driver~18 at Silverstone. The driver-preference setup bias is therefore not a source of cross-circuit error in the polar.

The inter-driver spread is reported as a diagnostic rather than a correction: at Jeddah, driver~18 yields $\hat\alpha = 0.968$ from three segments, a $10.4\%$ spread relative to driver~14's $0.876$. This spread likely reflects genuine intra-team setup differences (wing angle asymmetry, balance preferences, mechanical setup) rather than measurement noise alone. Because driver~14 is held fixed across all polar circuits, any systematic offset from driver-14-specific setup preferences is common-mode and cancels in the polar slope. However, the 10.4\% inter-car spread at Jeddah bounds the absolute accuracy of the $C_DA$ values as a proxy for the team's representative aerodynamic configuration; single-car absolute values may differ from the team's ``neutral'' reference by up to half this spread (approximately $\pm$5\%). If driver~14's preferred aero balance deviates systematically from the team's ``neutral'' reference setup, the absolute $C_DA$ and $C_LA$ values may be biased relative to what the other car would show.